\scanpage{17}
\mspage{16}

\uline{Ex} 2\qquad $X$ un schéma abélien $A$. On a ici $\operatorname{Todd}(X)=1$.

On trouve

\[
\chi_F(\delta)=\operatorname{tr}(\operatorname{ch}(L_\delta F))=\deg a_n
\]

si on pose

\[
\operatorname{ch}(L_\delta F)=\sum_1^n a_i\qquad a_i\in\operatorname{Chow}(A)
\]

On trouve ensuite, comme avant

\[
\varphi_F(\delta)=\inserted{\text{terme de degré 1 de }}(\exp D)\,.(a_{n-1})=\operatorname{pr}_{2*}(\tfrac12\underline{D^2a_{n-1}}).
\]

f)\qquad Nous allons donner une autre interprétation de l'homomorphisme $l_\delta$ de d), en définissant (en intégrant exclusivement sur $X\times_S B$) un homomorphisme

\[
\boxed{l^{\unsure{*|' }}\colon X\times\underline{\mathrm{Pic}}_{X/B}\longrightarrow\underline{\mathrm{Pic}}_{B/S}}.
\]

\margin{(additif en $\underline{\mathrm{Pic}}_{X/S}$)}

Pour ceci, soit $g$ une section de $X$, et $\lambda$ une section de $\underline{\mathrm{Pic}}_{X/S}$, qui (grâce à $g$) est définie par un module inversible $L\inserted{{}=L(\lambda,g)}$ sur \deleted{\unsure{le}} $X$, \unsure{trivialisée|trivialisé} le long de $g$. On peut alors former sur $X\times_S B$ ($W_g\inserted{{}=W(g)}$ faisceau de Weil sur $X\times_S B\unsure{)|^{g}}$, correspondant à $\unsure{\theta|0|g}$) $\det{}^{*}(\operatorname{pr}_{2*}(\operatorname{pr}_1^{*}(L)\otimes W))$ \deleted{\unsure{classe}} faisceau inv. sur $B$, dont la classe \deleted{\unsure{a}}\inserted{définit} une section $l(g,\lambda)$ de $\underline{\mathrm{Pic}}_{B/S}$. La formation est compatible avec extension de la base, d'où l'hom $l$.

Désignons, pour une section $g$ de $X$, par $l_g$ l'hom $\underline{\mathrm{Pic}}_{X/S}\mathbin{\unsure{-|\to}}\underline{\mathrm{Pic}}_{B/S}$ qu'il définit.
